% -------------------------------------------------------------------------
% WBS ID: RES-VER-CNV
% Deliverable: Convergence and Numerical-Error Test Methodology
% Status: Not started
% Evidence status: Not assessed
% Review status: Not reviewed
% Last updated:
% Dependencies:
% Downstream interfaces:
% -------------------------------------------------------------------------

\section{RES-VER-CNV --- Convergence and Numerical-Error Test Methodology}
\label{sec:res-ver-cnv}

\subsection{Purpose and Scope}
\label{sec:res-ver-cnv-purpose}

This Deliverable defines how the project separates numerical
convergence evidence from physical validation evidence. For the
Regional2D solver it covers spatial and temporal discretisation tests,
controlled manufactured-solution benchmarks, event-level
Tōhoku--Kamaishi forcing comparisons, and the interpretation of
mesh-to-mesh differences in exported coupling quantities. It does not
use observations, tune Manning friction, alter earthquake-source
amplitude, change bathymetric datasets, or calibrate the physical model.

The governing Regional2D mathematical model remains the accepted
nonlinear shallow-water system with conserved state
$\mathbf{U}=[h,q_x,q_y]^{\mathsf{T}}$, bed-slope physics, Manning
bottom friction, passive coseismic free-surface transfer, accepted
open/damped artificial boundary roles, wet/dry capability, and one-way
Regional2D-to-Local3D coupling quantities. The distinction between this
mathematical model and its numerical discretisation is mandatory:
first-order and limited-linear reconstruction are alternative numerical
methods for the same model, not different physical models.

\subsection{Research Questions}
\label{sec:res-ver-cnv-research-questions}

The Regional2D convergence programme asks:

\begin{enumerate}
  \item whether the implemented finite-volume solver is a consistent
    approximation of the accepted NLSWE model;
  \item whether event-specific forcing histories are stable under mesh
    refinement;
  \item whether mesh-to-mesh differences arise primarily from physical
    model fidelity, mesh geometry, source/terrain projection, temporal
    error, coupling extraction, or spatial discretisation error;
  \item whether a verified higher-order reconstruction reduces the
    event-level amplitude and phase sensitivity seen in the original
    first-order Tōhoku family;
  \item whether a spatially qualified Regional2D configuration can be
    frozen before any observational comparison or calibration.
\end{enumerate}

\subsection{Terminology and Definitions}
\label{sec:res-ver-cnv-definitions}

\begin{description}
  \item[Mathematical model] The continuous NLSWE, source, friction,
    terrain, boundary-role and coupling-variable specification.
  \item[Numerical discretisation] The finite-volume mesh, face-state
    reconstruction, hydrostatic reconstruction, numerical flux, SSPRK3
    time integration, CFL/timestep policy and mesh-size sequence.
  \item[Global solution convergence] The measured decrease in complete
    evolved-solution error under mesh refinement.
  \item[Local semi-discrete operator order] The order of a residual or
    truncation benchmark at fixed exact state. It is diagnostic evidence
    but is not by itself proof of global inconsistency.
  \item[Spatial qualification] A decision that the event-level forcing
    outputs are sufficiently insensitive to the selected mesh refinement
    sequence for the project tolerance and intended downstream use.
\end{description}

\subsection{Literature Evidence}
\label{sec:res-ver-cnv-literature-evidence}

% TODO: Record evidence from peer-reviewed papers, authoritative datasets,
% standards, technical reports and primary documentation.
%
% Do not create a detached literature-summary list. Explain how each source
% supports, contradicts or limits a project decision.

\subsection{Governing Relations and Quantitative Definitions}
\label{sec:res-ver-cnv-governing-relations}

The primary event forcing window is fixed as:

\begin{align}
  245
  \le
  t
  \le
  545~\mathrm{s}
  \label{eq:res-ver-cnv-r10-forcing-window}
\end{align}

The primary coupling quantities are the free surface $\eta(s,t)$, the
section-normal discharge per unit width $q_n(s,t)$, the integrated
normal discharge $Q_n(t)$ and the section-mean value $\bar{q}_n(t)$:

\begin{align}
  Q_n(t)
  &=
  \int_{\Gamma_c}
  q_n(s,t)
  \,\dd s
  \label{eq:res-ver-cnv-integrated-discharge}
  \\
  \bar{q}_n(t)
  &=
  \frac{Q_n(t)}{W}
  \label{eq:res-ver-cnv-mean-normal-discharge}
\end{align}

Distributed relative errors for two mesh levels $a$ and $b$ use the
established common-support definitions:

\begin{align}
  E_{\eta}
  &=
  \left[
    \frac{
      \int\!\!\int
      \left(
        \eta_a-\eta_b
      \right)^2
      \,\dd s\,\dd t
    }{
      \int\!\!\int
      \eta_b^2
      \,\dd s\,\dd t
    }
  \right]^{1/2}
  \label{eq:res-ver-cnv-distributed-eta-error}
  \\
  E_q
  &=
  \left[
    \frac{
      \int\!\!\int
      \left(
        q_{n,a}-q_{n,b}
      \right)^2
      \,\dd s\,\dd t
    }{
      \int\!\!\int
      q_{n,b}^2
      \,\dd s\,\dd t
    }
  \right]^{1/2}
  \label{eq:res-ver-cnv-distributed-qn-error}
\end{align}

\subsection{Test Objective}
\label{sec:res-ver-cnv-test-objective}

The C1A Regional2D objective is to establish a numerical baseline before
using any observational validation data. The original production method
was a cell-centred finite-volume scheme with piecewise-constant
hydrodynamic reconstruction, hydrostatic well-balanced reconstruction,
Rusanov flux, conservative face-based residual assembly and SSPRK3
hydrodynamic time integration. C1A-R8 verified this original method as a
globally first-order approximation of the accepted NLSWE model, with
finest-pair L2 orders approximately $0.9038$ for mass, $0.9129$ for
$q_x$ and $0.9574$ for $q_y$.

The event-level first-order Tōhoku family did not enter an asymptotic
forcing-convergence regime. With frozen terrain, source, physics,
domain and coupling extraction, the best first-order tested actual
characteristic mesh was approximately $260.45~\mathrm{m}$ for requested
h400. Adjacent R5 $Q_n$ waveform NRMSE values were approximately
$26.4\%$ for h600 to h500, $19.7\%$ for h500 to h450, and $21.6\%$ for
h450 to h400. This sequence was non-monotonic and must not be described
as mesh converged.

The R4D/R5 diagnosis excluded earthquake-source projection,
coupling-section extraction, independent $Q_n$ integration,
common-support interpolation, mass conservation, simple temporal phase
shift and poor h400 geometric mesh quality as primary causes. The
retained diagnosis was mixed spatial error: finite-volume
diffusion/dispersion as the primary mechanism and mesh-dependent
projection of the frozen bathymetry as a secondary mechanism. R6 mesh
quality supported this interpretation: minimum-angle p01 values were
approximately $44.18^\circ$ for h600, $45.51^\circ$ for h500,
$43.59^\circ$ for h450 and $45.88^\circ$ for h400, with no material
h400 radius-ratio degradation. The frozen processed terrain resolution
was $1000~\mathrm{m}$ while the solver actual characteristic sizes
reached approximately $388$, $325$, $296$ and $260~\mathrm{m}$; solver
mesh refinement therefore interpolated a fixed terrain representation
and was not equivalent to refining the physical terrain data.

The model/implementation traceability audit concluded
\texttt{MODEL\_CONSISTENT\_WITH\_DOCUMENTATION\_FIXES}, with
\texttt{MISMATCH = 0} and
\texttt{NOT\_IMPLEMENTED\_REQUIRED = 0}. This cleared the mathematical
model decision and left numerical discretisation as the appropriate
target for convergence improvement.

The R7 exact semi-discrete benchmark initially showed sub-first-order
local operator/truncation behaviour, but R8 separated local
semi-discrete operator order from global evolving-solution convergence.
R8 found a geometric closure residual of $0.0$, maximum first-moment
geometry residual of about $2.84\times10^{-14}$, face-orientation error
of $0.0$, and exact-face-quadrature final L2 errors of about
$9.07\times10^{-15}$ for mass, $3.59\times10^{-13}$ for $q_x$ and
$3.68\times10^{-13}$ for $q_y$. These results excluded mesh geometry
and exact physical-flux integration as fundamental defects and opened
the gate to investigate a higher-order reconstruction.

C1A-R9 introduced the opt-in \texttt{limited\_linear} reconstruction
while preserving \texttt{first\_order} as the default/reference path and
retaining \texttt{unlimited\_linear\_for\_verification} for controlled
checks. The method reconstructs $\eta$, $u$ and $v$ using weighted
least-squares gradients and a Barth--Jespersen limiter, leaves bed
states first order inside the hydrostatic reconstruction, and preserves
the Rusanov flux, SSPRK3, source physics, Manning friction, boundary
roles, wet/dry handling, coupling quantities and OpenMP residual
assembly. The controlled MMS finest-pair L2 orders were $2.3226$ for
mass, $2.2164$ for $q_x$ and $2.1614$ for $q_y$, giving classification
\texttt{SECOND\_ORDER\_VERIFIED}. At equal n32 resolution the L2 error
was reduced by approximately $30.37\times$, $8.63\times$ and
$6.05\times$ for mass, $q_x$ and $q_y$. A controlled smooth-wave proxy
reduced amplitude error by $5.35\times$ and phase-proxy error by
$17.43\times$. The measured medium-level runtime overhead was
approximately $1.223$, or about $22.3\%$, for substantially improved
controlled-problem accuracy.

The R10 event-level hypothesis is therefore precise: the
limited-linear reconstruction may reduce the first-order method's
Tōhoku--Kamaishi amplitude diffusion, phase dispersion and waveform
morphology sensitivity, but that must be measured on the frozen
h600/h500/h400 event mesh family before spatial qualification or any
temporal-convergence gate can be opened.

\subsubsection{Numerical and Physical Sensitivity Requirements}
\label{sec:res-ver-cnv-local-sensitivity-requirements}

The local-model validation programme shall distinguish:

\begin{enumerate}
  \item spatial and temporal discretisation sensitivity;
  \item incident-wave and boundary-condition sensitivity;
  \item turbulence-closure sensitivity;
  \item wall-treatment sensitivity;
  \item free-surface-capturing sensitivity;
  \item pressure-filtering and load-integration sensitivity;
  \item fluid-to-structure mapping sensitivity;
  \item structural time-integration and damping sensitivity.
\end{enumerate}

Grid refinement shall not be treated as sufficient evidence of model
validity. A refined solution may converge towards an incorrect result
where the boundary conditions, closure, constitutive model, or incident
state are inappropriate.

\subsection{Reference or Benchmark Solution}
\label{sec:res-ver-cnv-reference-benchmark-solution}

% TODO: Complete this RES-VER Domain-specific subsection.

\subsection{Test Configuration}
\label{sec:res-ver-cnv-test-configuration}

% TODO: Complete this RES-VER Domain-specific subsection.

\subsubsection{Spatial Resolution Requirements}
\label{sec:res-ver-cnv-spatial-resolution-requirements}

Grid adequacy shall be determined for each output separately.

A grid that reproduces a deep-ocean leading-wave amplitude shall not be
assumed adequate for:

\begin{itemize}
  \item narrow bays;
  \item steep ria coastlines;
  \item local run-up;
  \item river propagation;
  \item coastal structures;
  \item barrier-scale pressure and force.
\end{itemize}

Computational-grid spacing and source-data resolution shall be reported
independently. Refinement of the computational mesh below the
resolution of the bathymetric or topographic input shall not be treated
as equivalent physical refinement.

\subsection{Measured Quantities}
\label{sec:res-ver-cnv-measured-quantities}

% TODO: Complete this RES-VER Domain-specific subsection.

\subsection{Error Metrics}
\label{sec:res-ver-cnv-error-metrics}

For each event mesh pair, the analysis shall report RMSE, NRMSE,
correlation, peak value, peak time and peak-time difference for
$\eta$, $q_n$, $Q_n$ and $\bar{q}_n$. Formal convergence remains based
on physical-time waveforms; diagnostic phase alignment may be reported
separately through physical-time NRMSE, optimal lag, phase-aligned NRMSE
and phase contribution.

First-order histories are historical same-mesh comparisons, not
references for the limited-linear convergence estimate. The required
method-impact comparison is first-order versus limited-linear at the
same h600, h500 and h400 meshes, using peak amplitudes, phase,
waveform morphology and, where useful, spectral content.

\subsection{Tolerance and Acceptance Criteria}
\label{sec:res-ver-cnv-tolerance-acceptance-criteria}

The project target for medium-to-fine principal forcing differences is
approximately $2\%$. For R10 the primary spatial pair is h500 to h400,
because h400 is the fine member of the minimum-cost limited-linear
event family. Spatial qualification requires principal h500--h400
forcing differences near or below this threshold, no material waveform
morphology change, no large unresolved physical-time phase difference,
stable distributed forcing, and passing terrain/source/domain
invariance. Qualification shall not be granted from peak $\eta$ alone.

If h500--h400 is slightly above $2\%$ but monotonic, strongly reduced
relative to first order, and clearly approaching an asymptotic regime,
the result shall be classified as
\texttt{APPROACHING\_SPATIAL\_QUALIFICATION}, not qualified. Strongly
oscillatory behaviour shall be classified as
\texttt{SPATIAL\_NON\_ASYMPTOTIC}. Temporal convergence may proceed
only after \texttt{SPATIAL\_QUALIFIED}.

\subsection{Execution Handoff}
\label{sec:res-ver-cnv-execution-handoff}

R10 shall use the exact frozen R4/R5 physical family: conditioned
terrain SHA-256
\texttt{45e5ab63a69e77ec11b293c39cbb93dd0df30a4f24d1a4f4d9515267a01f1363},
source SHA-256
\texttt{88f58bb256c8e5ff7baa8ec662118572b1e6a1b38b0fdd9b85a0541ddf6f6498},
the same domain, Manning coefficient, gravity, boundary policy, sponge
policy, wet/dry settings, coupling section, output cadence and
$600~\mathrm{s}$ horizon. The required meshes are the historical h600,
h500 and h400 artefacts where available, enabling direct first-order
versus limited-linear same-mesh comparison. Requested h450 is excluded
from the minimum R10 family because its previous actual refinement ratio
provided weak convergence leverage for substantial additional compute.

The R10 production candidate must select
\texttt{reconstruction = limited\_linear} explicitly. The first-order
method remains available as a backend and historical comparator.

\subsection{Candidate Approaches}
\label{sec:res-ver-cnv-candidate-approaches}

% TODO: Define the credible candidate approaches considered.

\subsection{Critical Comparison}
\label{sec:res-ver-cnv-critical-comparison}

% TODO: Compare candidates using physically and project-relevant criteria.
% Do not select an approach solely on popularity or implementation ease.

\subsection{Assumptions and Applicability}
\label{sec:res-ver-cnv-assumptions}

% TODO: State assumptions and the conditions under which they are valid.

\subsection{Limitations and Failure Conditions}
\label{sec:res-ver-cnv-limitations}

% TODO: State where the evidence, model or selected approach ceases to be
% reliable or applicable.

\subsection{Project-Specific Conclusions}
\label{sec:res-ver-cnv-conclusions}

Entering R10, the project conclusion is that the original Regional2D
implementation was a verified globally first-order finite-volume
approximation of the accepted NLSWE mathematical model, but it exhibited
large non-asymptotic forcing sensitivity in the Tōhoku--Kamaishi
corridor. A limited-linear reconstruction was scientifically justified
as a numerical-discretisation improvement, not as a physical-model
change. Controlled C1A-R9 MMS and smooth-wave tests verified
approximately second-order behaviour and substantially reduced
amplitude and phase error. The remaining numerical question is whether
that controlled improvement restores practical spatial convergence for
the event-specific Tōhoku--Kamaishi forcing.

\subsection{Selected Approach and Rationale}
\label{sec:res-ver-cnv-selected-approach}

% TODO: Record the selected approach only after sufficient evidence exists.
% State the alternatives rejected and the reasons for rejection.

\subsection{Unresolved Decisions}
\label{sec:res-ver-cnv-unresolved-decisions}

% TODO: Record unresolved questions, required evidence and decision owners.

\subsection{Requirements and Handoffs}
\label{sec:res-ver-cnv-requirements}

The Software C1A evidence tree is the authority for detailed numerical
metrics:
\texttt{docs/workstream/SWE - Software/SWE-VER/SWE-VER-CONV/C1A/}.
R10 shall add event evidence there, keep raw outputs under an external
R10 root, and update this Research workstream only after actual event
results are available. No observational comparison, calibration,
Local3D run or h300 run may proceed from this Deliverable unless the
Regional2D numerical gates explicitly allow it.

\subsection{Deliverable Completion Record}
\label{sec:res-ver-cnv-completion-record}

% TODO: Record whether the Deliverable is:
%
% - Not started
% - In progress
% - Draft complete
% - Reviewed
% - Accepted
%
% Acceptance requires:
%
% - the research questions to be answered;
% - claims to be supported by appropriate evidence;
% - assumptions and limitations to be explicit;
% - the project-specific conclusion to be stated;
% - downstream requirements to be traceable;
% - unresolved decisions to be closed or formally deferred.
